\subsection{Landscape, the Scale of SUSY Breaking, and Inflation --- arXiv:hep-th/0411011}

\textbf{Authors:} Renata Kallosh, Andrei Linde

\paragraph{主な主張}
最も単純なKKLT模型ではインフレーション中の安定化から$H\lesssim m_{3/2}$が要求される一方、複数の指数項を持つracetrack型模型ではこの制約を回避できることを示した\cite{Kallosh:2004yh}。

\paragraph{新規性}
体積モジュラスをSUSY Minkowski真空に安定化する構成によって障壁の高さを小さいgravitino質量から切り離し、高いインフレーションスケールとの両立を可能にした。

\paragraph{位置付け}
インフレーションとSUSY破れスケールの関係を明確にした、KL型racetrack安定化の基礎文献。

\paragraph{コメント（読書メモ）}
元メモでは「racetrack model元論文」と記録。ここでは、racetrack超ポテンシャルを用いて軽いgravitinoと高いインフレーションスケールを両立するKL模型の原論文として参照する。
